Um zu zeigen, dass $\epsilon^{\prime \alpha \beta} = -\epsilon^{\alpha \beta}$, betrachten wir die Definition der inversen Transformation. Die Levi-Civita-Symbole $\epsilon^{\alpha \beta}$ und $\epsilon^{\prime \alpha \beta}$ sind antisymmetrische Tensoren, die in zwei Dimensionen definiert sind. Die Transformation eines solchen Tensors unter einer linearen Transformation $A$ ist gegeben durch:

\[
\epsilon^{\prime \alpha \beta} = A^\alpha_{\ \gamma} A^\beta_{\ \delta} \epsilon^{\gamma \delta}
\]

Da $\epsilon^{\alpha \beta}$ antisymmetrisch ist, gilt $\epsilon^{\alpha \beta} = -\epsilon^{\beta \alpha}$. Für die inverse Transformation $A^{-1}$ gilt:

\[
\epsilon^{\alpha \beta} = (A^{-1})^\alpha_{\ \gamma} (A^{-1})^\beta_{\ \delta} \epsilon^{\prime \gamma \delta}
\]

Da $A$ eine orthogonale Transformation ist, gilt $A^{-1} = A^T$, und somit:

\[
\epsilon^{\prime \alpha \beta} = A^\alpha_{\ \gamma} A^\beta_{\ \delta} \epsilon^{\gamma \delta} = -\epsilon^{\alpha \beta}
\]

Dies zeigt, dass die transformierte Version des Levi-Civita-Symbols das negative des ursprünglichen Symbols ist, was die geforderte Beziehung bestätigt.